% Research notes for the parent article. Classical results are explicitly
% attributed; the sharp suffix-specific three-jump bound is proved below.
\section{Weighted transport on a checked orbit trace}

Let $I_N=\{0,\ldots,N-1\}$ and let $\sim$ be the equivalence relation
generated by the input interval isometries. The input weight function is
$w:I_N\to\mathbb Z^d$, represented by an additive list of $m$ half-open
constant intervals. An endpoint sweep converts it into a canonical complete
partition with at most $2m+1$ runs, including any zero-weight runs. For a
component $O$ the requested weight is
\[
             W(O)=\sum_{x\in O}w(x).
\]
The output is a finite list of distinct pairs $(z,\mu_z)$, where $\mu_z$
counts the components whose weight vector equals $z$. This represents all
components without enumerating them. Zero vectors, negative coordinates,
and $d=0$ are legitimate; zero weight never means that a component is absent.

This is an implementation of the weighted Agol--Hass--Thurston algorithm
\cite[Section 6]{aht2006},
not a new general orbit-counting theorem. The existing scheduler supplies
a complete local orbit proof. The ordinary independent verifier checks
that proof, after which a weight observer follows the same operations.
The observer does not make scheduling decisions. Identity deletion,
reflection trimming, the sharp or classical periodic merger, and
transmission preserve the actual equivalence relation on the current
points, so they leave their weights unchanged.

\subsection{The exact full-range transfer}

Immediately before a truncation the carrier has range $[c,N)$ and domain
$[a,b+1)$, with $a<c$. The entire range is transferred, including its part
that will remain after truncation. This full transfer is essential for the
small breakpoint bound; repeatedly transferring only removed pieces is
not the operation analysed here.

For a positive carrier put $p=c-a$. A range point $x$ is moved to
\[
             \rho(x)=a+((x-a)\bmod p)\in[a,c).
\]
This is obtained by iterating the carrier's partial inverse until its image
first lies before $c$. Every intermediate inverse is defined. If the
domain and range are disjoint, only one inverse is needed, and the same
formula applies.

Suppose $w$ equals $v$ on the range subinterval $[l,h)$. Write
$h-l=qp+r$, $0\le r<p$. The pushforward of this interval adds $qv$ on
all of $[a,c)$ and one more copy of $v$ on the cyclic interval of length
$r$ beginning at $a+((l-a)\bmod p)$. A cyclic interval is represented by
one ordinary interval if it does not wrap and by two if it does. The
formula uses one integer division; it performs no loop of length $q$.

For a trimmed reversing carrier the range point $x$ is moved once to
$a+N-1-x$. A constant range interval $[l,h)$ therefore adds its value to
$[a+N-h,a+N-l)$. In either case the old range weights are first set to
zero. Only their transferred copies remain. Every original contribution
is moved to exactly one point in the same orbit, proving preservation of
every orbit sum, even for signed vector weights. The existing truncation
rule then removes a zero-weight suffix whose points have earlier orbit
representatives.

At static contraction, every removed point represents one complete orbit.
Intersect each static gap with the current constant-weight runs and add
its integer length to the multiplicity of that vector. Shift the remaining
run endpoints by deleted rank, just as the ordinary interval points are
shifted. This gives two executable invariants:
\[
 \sum_z\mu_z=C,\qquad
 \sum_z\mu_z z=\sum_{x\in I_N}w(x),
\]
where $C$ is the independently verified orbit count. These invariants are
necessary checks; they alone would not prove the individual component
weights, which are justified by the local transport argument.

\subsection{A sharp suffix-specific run bound}

The general AHT transfer lemma permits an increase of four constant
intervals. Because the maintained scheduler always transfers a full
rightmost range, its bound improves to three.

\begin{proposition}[Three-run suffix bound]
Suppose the current nonempty universe has $m$ constant-weight runs.
Transfer the full carrier range $[c,N)$ as above and truncate to
$I_{N'}$, where $c\le N'<N$. The resulting weight function has at most
$m+3$ runs. If $N'=c$, or if $a=0$, it has at most $m+2$ runs.
The bound $m+3$ is attained by a legal step of the maintained scheduler.
\end{proposition}

\begin{proof}
Let $J=\{j:0<j<N,\ w(j-1)\ne w(j)\}$, so $m=1+|J|$.
Put $s=|J\cap(c,N)|$. Old changes strictly below $c$ may remain. The
$s$ old changes inside the transferred range disappear there.

For a disjoint positive carrier or a trimmed reflection, a pushforward
has changes only at the two source endpoints and at the images of those
$s$ old changes. Setting the surviving part of $[c,N)$ to zero may add
a change at $c$. Thus there are at most $s+3$ new potential changes,
which replace the $s$ removed changes.

For an overlapping positive carrier, regard the pushforward on
$[a,c)$ as a cyclic residue sum. Its interior changes can occur only at
residues of the $s$ old changes and at the residue of $N$. The initial
range endpoint $c$ has residue $a$. Extending the pushforward by zero
can add source-boundary changes at $a$ and $c$. Hence all changes again
belong to the old changes below $c$, the at most $s$ folded old changes,
and three additional points: $a$, $c$, and the residue of $N$.
Coincident points, zero vector changes and pre-existing changes only
decrease their number. A source endpoint at zero is not an interior
change. Nor is $c$ an interior change when $N'=c$, giving the stated
improvements.

For sharpness, let $N=7$, use constant weight $1$, and take the three
pairings
\[
 [1,2]\longrightarrow[5,6],\qquad
 [3,4]\longrightarrow[4,5],\qquad
 [0,0]\longrightarrow[3,3].
\]
All points are initially covered. The first pairing is the rightmost
carrier. There is no periodic merger and neither remaining range lies
inside its range. The second pairing forces the truncation point
$N'=6$. Full-range transfer produces values $1,2,1,0$ respectively on
$[0,1),[1,3),[3,5),[5,6)$. Thus one run becomes four.
\end{proof}

\begin{figure}[htbp]
\centering
\includegraphics[width=0.84\textwidth]{sharp_transfer.pdf}
\caption{The exact seven-point example in the proof. Transferring the carrier
range $[5,7)$ to $[1,3)$ and removing point $6$ increases the number of
constant runs from one to four. Zero weight at surviving point $5$ remains
part of the representation.}
\label{fig:sharp-transfer}
\end{figure}

Static contraction cannot increase the number of weight runs: it merely
takes an ordered subsequence of point values and coalesces equal adjacent
runs. All other operations leave weights fixed. Consequently, after $T$
full transfers the current weight partition has at most $m_0+3T$ runs.
This bound is independent of the magnitudes of all represented intervals
and multiplicities. It remains valid with the Fine--Wilf merger because
that merger does not alter point weights or the equivalence relation.

For signed input weights let
$A=\sum_j (h_j-l_j)\|v_j\|_1$, using the additive input intervals.
Every intermediate coordinate has absolute value at most $A$: weights
are gathered, never copied into two different orbit representatives.
The bit length is therefore $O(1+\log(A+1))$, bounded by
$O(\log(N+1)+\log(m+1)+L+\log(d+1))$ when input coordinate bit length
is at most $L$. The degree-zero case is treated separately with empty
vectors. With the existing polynomial-size orbit trace, these facts give
polynomial bit complexity for canonicalization, all vector arithmetic,
and the compressed output. No new polynomial degree for the scheduler is
asserted. A simple bound on the number of emitted distinct records is the
number of trace contractions times $m_0+3T$, still polynomial.

\subsection{Sparse marked-component incidence}

For each of $r$ ports $A_i$ represented as a union of intervals, assign
coordinate $i$ of $w(x)$ to be $1$ if $x\in A_i$ and $0$ otherwise.
Normalize the intervals within each port so overlaps do not accidentally
count a point twice. The weighted output gives
\[
        W_i(O)=|O\cap A_i|,
\]
which contains both exact point multiplicity and the incidence signature
$\{i:W_i(O)>0\}$. Merge output vectors having the same support to obtain
a sparse histogram of attained signatures. Empty support is retained.
With $M$ marked intervals, this computation has polynomial bit complexity
in $k,M,r,\log(N+1)$; there is no $2^r$ output requirement. If the caller
requests a dense array indexed by every subset of $[r]$, that conversion
still requires $\Omega(2^r)$ slots. This sparse interface is an application
of the classical weighted theorem, not a new general result about AHT.

Additive incidence does not determine attachment order. A single circle
with four distinctly labelled marked points in cyclic order $a,b,c,d$
has the same weight vector $(1,1,1,1)$ as one with order $a,c,b,d$.
These orders are inequivalent under rotation and reversal with the labels
fixed. Thus even exact multiplicities and supports cannot recover the
ordered boundary data needed by a general cutting or gluing procedure.
Such a procedure must retain an additional order and transport structure.

\subsection{Connected normal-coordinate extraction}

For each of the at most $7t$ normal-disc types choose one incident local
edge. Add the corresponding basis vector at precisely the edge points
met by that disc stack on that chosen local edge. A triangle at vertex
$v$ occupies the points nearest $v$. On an edge directed $a$ to $b$, an
incident quadrilateral stack occupies the interval after the $a$-triangle
stack and before the $b$-triangle stack. The actual global edge orientation
may reverse the interval; obtain it from the validated parity data.
Different local edges may represent one global edge, so these basis
contributions are additive and must not overwrite each other.

Each local normal disc contributes one basis unit exactly once. Therefore
the orbit weight is the complete normal-coordinate vector of that actual
connected component. Grouped multiplicities encode equal components.
Their coordinate sum with multiplicity recovers the supplied vector.
This construction works for one-sided components directly and needs no
common-divisor reduction. In particular, $2$ times a one-sided component
may describe a connected two-sided double instead of two components;
the orbit extraction, rather than linear multiplicity guessing, decides
which occurs. This is the classical AHT Corollary 17 construction made
executable over the maintained edge-arc encoding.

Once weighted extraction has proved connectedness, Euler characteristic
one and nonempty boundary suffice to identify a disc, as proved in
Lemma~\ref{lem:disk-shortcut}. The existing torus-boundary homology test
then decides whether its boundary is essential. No further boundary or
orientation orbit query is needed. Extracting such a vector from
a disconnected input is new functionality of this code base. It does not
provide surface search, diagram-to-exterior provenance, hierarchy
construction, or a quasi-polynomial complete recognition procedure.

\subsection{Exact scope of replay and limits}

The weight certificate binds the normalized weight function and dimension,
the source-bound unweighted trace, and the ordered exact output classes.
Replay invokes no orbit scheduler. The ordinary orbit checker is
independent of orbit search. The weight producer and weight checker share
the small transport arithmetic; its independent validation comes from
the expanded graph and direct pushforward tests. Calling the entire
weight verifier arithmetically independent of weight production would
overstate this boundary.

The producer's cycle limit applies only to ordinary orbit search. The
weight-run, output-record and event limits are checked in the subsequent
weighted replay; the full ordinary trace has already been generated and
validated by then. They are cooperative work/output limits, not memory
or wall-clock guarantees. The standalone replay and weight-certificate
checker reject an over-limit trace by its explicit operation count before
ordinary proof replay. Every incomplete producer or standalone replay
returns no orbit count, class list, aggregate or certificate. Cancellation
exceptions propagate.
